% ch10：§9 附录（原文 p59–p63，至 References 前）
\section{附录：复情形下的构造}

在上一节中我们考察了完备局部非阿基米德域 $\C_{\eps}$ 上的代数簇与解析簇。本节说明在复数情形下（即假定 $\eps$ 是一个固定正数）我们的进路。应当提醒读者：在这一情形下如何获得同调镜像猜想的严格证明，目前尚不清楚。特别地，人们还不知道如何证明定义 Fukaya 范畴中复合的那些级数的收敛性。尽管如此，我们仍要讨论复情形，因为这里的几何更透明。例如，可以用实形式的 Fourier--Mukai 变换在对象上构造镜像对称函子（见 \S9.1）。从本文主体的观点看，附录可以看作一个几何上的动机。出于这个理由，我们不再强调 $X$ 是阿贝尔簇，而将使用我们关于坍缩的猜想，以及「环面纤维化的底 $Y$ 是带整仿射结构与凯勒势的光滑流形」这一假设。我们也将沿用 \S3 的记号。

\subsection{复情形下对象上的镜像对称函子}

在复数情形，镜像对称函子把一对 $(L,\rho)$ 对应到 $X=X_{\eps}$ 上的全纯向量丛 $F(L,\rho)$，其中 $L\subset X^{\vee}$ 是拉格朗日子流形，使得投影 $p^{\vee}_L:L\to Y$ 是非分歧覆盖，而 $\rho$ 是 $L$ 上的局部系。若 $L$ 是 $p^{\vee}$ 的一个截面且 $\rank(\rho)=1$，则 $E=F(L,\rho)$ 是线丛。一般地，$E$ 可以局部表示为直和 $E\simeq\oplus_{\alpha\in A}E_{\alpha}$，其中 $A$ 是覆盖 $L\to Y$ 的叶（即连通分支）的集合，而 $E_{\alpha}$ 是秩等于 $\rho$ 在叶 $\alpha$ 上的秩的全纯向量丛。

对象上镜像对称函子的下述显式构造并非新的，见例如 \cite{AP}。我们从一个注记开始：在 $X\times_Y X^{\vee}$ 上存在典范的 $U(1)$-丛（Poincaré 线丛）。记之为 $P$。它带有典范联络，其描述见下。固定 $y\in Y$。则 $p^{-1}(y)\simeq T_{Y,y}/\eps T^{\Z}_{Y,y}$，而 $(p^{\vee})^{-1}(y)\simeq T^{*}_{Y,y}/(T^{\Z}_{Y,y})^{\vee}$。我们把环面 $(p^{\vee})^{-1}(y)$ 与「环面 $p^{-1}(y)$ 上、在点 $0\in p^{-1}(y)$ 处平凡化的 $U(1)$-局部系的模空间」等同起来。我们定义 $U(1)$-丛 $P$ 为 $X\times_Y X^{\vee}$ 上对应于这一描述的万有丛。

为描述 $P$ 上的联络，考虑纤维方向的万有覆盖 $r:T_Y\to T_Y/T^{\Z}_Y$ 与 $r^{\vee}:T^{*}_Y\to T^{*}_Y/(T^{\Z}_Y)^{\vee}$。则 $P$ 拉回到 $T_Y\times_Y T^{*}_Y$ 上的丛 $\tilde P$ 可典范平凡化。于是可以在坐标中工作。设 $y=(y_1,\dots,y_n)$ 是 $Y$ 上的坐标，$x=(x_1,\dots,x_n)$ 与 $x^{\vee}=(x_1^{\vee},\dots,x_n^{\vee})$ 分别是 $T_Y\to Y$ 与 $T^{*}_Y\to Y$ 的纤维上的坐标。甲板变换 $x_j\mapsto x_j+\eps n_j$，$n_j\in\Z$，作用在 $\tilde P$ 上并保持平凡化，而变换 $x_j^{\vee}\mapsto x_j^{\vee}+n_j^{\vee}$，$n_j^{\vee}\in\Z$，则以乘子 $\exp\bigl((2\pi i/\eps)\sum_j n_j^{\vee}x_j\bigr)$ 作用在 $\tilde P$ 上。

设 $\nabla_0$ 是 $\tilde P$ 上的平凡联络。考虑 $\tilde P$ 上由下列公式给出的联络 $\bar\nabla$：
\[
\bar\nabla=\nabla_0+\frac{2\pi i}{\eps}\sum_{1\le j\le n}x_j^{\vee}\,dx_j.
\]

\begin{lmm}
联络 $\bar\nabla$ 给出 $P$ 上的一个联络。
\end{lmm}

\begin{proof}
显然，$\bar\nabla$ 在变换 $x_j\mapsto x_j+\eps n_j$，$n_j\in\Z$ 下不变。变换 $x_j^{\vee}\mapsto x_j^{\vee}+n_j^{\vee}$，$n_j^{\vee}\in\Z$ 连同 $\bar\nabla$ 由 $h=\exp\bigl((2\pi i/\eps)\sum_j n_j^{\vee}x_j\bigr)$ 给出的规范变换，也保持 $\bar\nabla$ 不变。引理得证。
\end{proof}

仍设 $(L,\rho)$ 如上。镜像对称函子把它对应到全纯向量丛 $E=F(L,\rho)$，使得（在坐标中）它在点 $(y,x)$ 上的纤维由公式
\[
E(y,x)=\bigoplus_{x^{\vee}\in L,\ p^{\vee}(x^{\vee})=y}\rho(x^{\vee})\otimes P(x,x^{\vee})
\]
给出。该向量丛带有诱导联络 $\nabla_E$。当 $\rho$ 是酉的时，$E$ 还带有自然的 Hermite 度量。

\begin{prp}
$\operatorname{curv}(\nabla_E)$ 的 $(0,2)$-部分是平凡的。特别地，$\nabla_E$ 是全纯联络。
\end{prp}

\begin{proof}
这可由 $L$ 是拉格朗日这一事实推出。事实上，把 $L$ 提升到 $T^{*}_Y$。则在 $L$ 的任一连通分支的局部邻域内，可以找到光滑实函数 $f=f(y)$ 使 $L=df$。于是 $L$ 的局部方程为 $x_j^{\vee}=\p f/\p y_j$，$1\le j\le n$。联络 $\nabla_E$ 局部可以写成
$\nabla_E=\nabla_{E,0}+i\,d_E\otimes\bigl((2\pi i/\eps)\sum_j (\p f/\p y_j)\,dx_j\bigr)$，
其中 $\nabla_{E,0}$ 是向量丛 $E$ 上的平凡平坦联络。由于 $T_Y$ 上的全纯坐标由 $z_j=y_j+i\,x_j$（$i=\sqrt{-1}$）给出，可见曲率的 $(0,2)$-部分等于
\[
\operatorname{curv}(\nabla_E)^{(0,2)}=\text{const}\times\Bigl(\sum_{j,k}\frac{\p^2 f}{\p y_j\,\p y_k}\,d\bar z_j\,d\bar z_k\Bigr)=0 .\qedhere
\]
\end{proof}

\begin{dfn}
对 $X$ 上任意两个全纯向量丛 $E_1$、$E_2$，定义
$\Hom_{\mathrm{Dolb}}(E_1,E_2)=\Omega^{0,*}\bigl(X,\Hom(E_1,E_2)\bigr)$。
\end{dfn}

我们把取值于向量丛 $\Hom(E_1,E_2)$ 的 Dolbeault 形式构成的空间看作关于 $\dbar$-微分的微分分次代数。这样就得到了 $X$ 上凝聚层的导出范畴上的 $\Ainf$-范畴结构（实际上是 dg-范畴结构）。可以证明，这一 $\Ainf$-结构与正文中提到的那个结构是等价的。

\subsection{Dolbeault 形式空间中的扇区}

设 $E_i=F(L_i,\rho_i)$，$i=1,2$ 如上。在复数情形，存在与非阿基米德几何中的 dg-范畴 $\mathcal C(Y)$ 类似的范畴，记作 $\A(Y)$。$\A(Y)$ 的对象是 $X$ 上形如 $E=F(L,\rho)$ 的全纯向量丛。态射是 $Y$ 上软层的截面。确切地说，定义 $Y$ 上的层 $\underline{\Hom}_{\A(Y)}(E_1,E_2)$ 为正向像 $p_*\bigl(\underline{\Hom}_{\mathrm{Dolb}}(E_1,E_2)\bigr)$（记号自明）。于是 $\Hom_{\A(Y)}(E_1,E_2)$ 是该层的整体截面。该层对应于非阿基米德几何中所考虑的层 $\underline{\Hom}_{\mathcal C(Y)}(E_1,E_2)$。取 $Y$ 的开仿射卡 $U\subset Y$，$U\simeq\R^n$。则 $\Gamma\bigl(U,\underline{\Hom}_{\A(Y)}(E_1,E_2)\bigr)$ 含有（相对于环面 $T^n$ 在 $\Gamma\bigl(U\times T^n,\underline{\Hom}_{\mathrm{Dolb}}(E_1,E_2)\bigr)$ 上的自然作用而言的）有限 Fourier 和构成的子层。这样便得到层 $\underline{\Hom}^{\mathrm{alg}}_{\A(Y)}(E_1,E_2)$，它是非阿基米德情形所考虑的层 $\underline{\Hom}^{\mathrm{alg}}_{\mathcal C(Y)}(E_1,E_2)$ 的类似物。注意存在层的自然同态 $j:p^{*}(\Omega^{0,*}_Y)\to\Omega^{0,*}_X$。$j$ 的像由沿 $p$ 的纤维为局部常系数的 Dolbeault 形式组成。在局部坐标中，$j$ 由公式
\[
f_{i_1\dots i_n}(y)\,dy_{i_1}\wedge\dots\wedge dy_{i_n}\longmapsto
f_{i_1\dots i_n}(y)\,d\bar z_{i_1}\wedge\dots\wedge d\bar z_{i_n}
\]
给出，其中 $\bar z_k=y_k-\sqrt{-1}\,x_k$，$1\le k\le n$。容易看出，$j$ 与 de Rham 形式、Dolbeault 形式上的 dg-代数结构相容。于是对 $X$ 上的一对全纯向量丛 $E_1,E_2$，层 $\underline{\Hom}_{\A(Y)}(E_1,E_2)$ 上带有 $\Omega^{0,*}_Y$ 上的典范 dg-模结构。当 $E_i=F(L_i,\rho_i)$，$i=1,2$ 时，子层 $\underline{\Hom}^{\mathrm{alg}}_{\A(Y)}(E_1,E_2)$ 也是 $\Omega^{0,*}_Y$ 上的 dg-模层。

与非阿基米德情形一样，茎 $\underline{\Hom}^{\mathrm{alg}}_{\A(Y)}(E_1,E_2)_y$，$y\in Y$，有典范分解，为 $\Omega^{0,*}_{Y,y}$ 上有限秩 dg-模的直和。被加项由同伦类 $[\gamma]\in P(L_1,L_2,y)$ 标号，称为扇区。把它们记作 $\underline{\Hom}^{\mathrm{alg},[\gamma]}_{\A(Y)}(E_1,E_2)_y$。非正式地说，扇区对应于 Dolbeault 形式沿环面纤维方向的「Fourier 分量」。下面更具体地描述它们。为简单起见，假定 $\rho_i$，$i=1,2$，是秩一的平凡局部系，且 $L_i$，$i=1,2$，与 $p$ 的每个纤维恰好交于一点。则在 $p^{-1}(y)$ 附近可以写 $L_i=\operatorname{graph}(df_i)\ (\mathrm{mod}\ (T^{\Z}_Y)^{\vee})$，$i=1,2$，其中 $f_i$ 是 $Y$ 上光滑函数在 $y$ 处的芽。由 Poincaré 丛的描述可以推出，$\underline{\Hom}^{\mathrm{alg}}_{\A(Y)}(E_1,E_2)_y$ 典范地等同于 $T^n=p^{-1}(y)$ 附近的 $\dbar$-形式的芽空间，带有扭转微分
\[
\dbar_{\alpha}\alpha=\dbar\alpha+\frac{i}{\eps}\sum_j\frac{\p f}{\p y_j}\,d\bar z_j\wedge\alpha,
\]
其中 $f=f_1-f_2$。于是对应于路径 $\gamma$ 的扇区由形如
\[
\alpha=\sum_{i_1,\dots,i_n}\exp\bigl(i\langle m,x\rangle/\eps\bigr)f_{i_1\dots i_n}(y)\,d\bar z_{i_1}\wedge\dots\wedge d\bar z_{i_n}
\]
的 Dolbeault 形式组成。此处向量 $m=m(\gamma)$ 是 $T^n=p^{-1}(y)$ 中一个闭路的同伦类，该闭路是如下三条路径的复合：
\begin{enumerate}[label=\arabic*),leftmargin=2.2em,itemsep=2pt,topsep=3pt]
\item 路径 $[0,1]\to T^n$，$t\mapsto t\,(df_1)_y\ \mathrm{mod}\ (T^{\Z}_Y)^{\vee}$；
\item 路径 $\gamma$；
\item 路径 $[0,1]\to T^n$，$t\mapsto(1-t)(df_2)_y\ \mathrm{mod}\ (T^{\Z}_Y)^{\vee}$。
\end{enumerate}

在非阿基米德情形，扇区的选取对应于单项式 $z_1^{i_1}\dots z_n^{i_n}$ 的选取。在非阿基米德进路中，路径的同伦类对应于 Fourier 级数的被加项。局部地，每个扇区可以等同于 $Y$ 上的 de Rham 复形。确切地说，对形式
$\alpha=\sum_{i_1,\dots,i_n}f_{i_1\dots i_n}(y)\exp(\langle m,x\rangle)d\bar z_{i_1}\wedge\dots\wedge d\bar z_{i_n}$，
我们对应形式
$\alpha_m=\sum_{i_1,\dots,i_n}f_{i_1\dots i_n}(y)\exp\bigl(\tfrac{1}{\eps}\langle m,x\rangle\bigr)dy_{i_1}\wedge\dots\wedge dy_{i_n}$，
其中 $m=m(\gamma)$ 界定了扇区。容易看出，$X$ 上 Dolbeault 形式上的微分 $\dbar$ 对应于 $\Omega^{*}(Y)$ 上的 de Rham 微分 $d$。这样便得到复形的同构
$\underline{\Hom}^{\mathrm{alg},[\gamma]}_{\A(Y)}(E_1,E_2)_y\simeq\Omega^{*}_{Y,y}\otimes\C$。

\begin{rmk}
当 $\eps$ 不是固定的数，而是参数 $\eps\to 0$ 时，系数 $f_{i_1\dots i_n}(y)$ 是 $\eps$ 的渐近级数，形如
\[
f_{i_1\dots i_n}(y,\eps)=\sum_{j\ge 1}\exp(-\lam_j/\eps)f_{j,i_1\dots i_n},
\]
其中 $\lam_j\in\R$，$\lam_1<\dots<\lam_j<\dots$，且 $\lam_j\to+\infty$。
\end{rmk}

出现在 $\alpha_m$ 于 $y$ 处展开式中的指数集合，对应于非阿基米德情形所考虑的谱 $\operatorname{Sp}_y(\alpha)$。

\subsection{半群 $\varphi^t$}

现在可以定义半群 $\varphi^t:\Hom_{\mathrm{Dolb}}(E_1,E_2)\to\Hom_{\mathrm{Dolb}}(E_1,E_2)$，$0\le t<+\infty$。它是非阿基米德情形中半群 $\varphi^t$ 的类似物。首先像上面那样，把扇区 $\underline{\Hom}^{\mathrm{alg},[\gamma]}_{\A(Y)}(E_1,E_2)_y$ 与 $\Omega^{*}\bigl(Y,\Hom(\rho_1,\rho_2)\bigr)_y$ 等同起来。回忆非阿基米德部分：与路径的同伦类 $\gamma$ 典范地相联系着一个闭 $1$-形式 $\mu_{\gamma}=\int_{\gamma}\omega$，其中 $\omega$ 是 $X^{\vee}$ 上的辛形式。利用 $Y$ 上的黎曼度量 $g_Y$，把 $\mu$ 对应到 $Y$ 上的向量场 $\xi_{\gamma}$。在局部平凡化中，它由 $\grad\bigl((f_1-f_2+\langle m(\gamma),\cdot\rangle)/\eps\bigr)$ 给出。于是 $\varphi^t$ 的无穷小作用定义为 Lie 导数 $\Lie_{\xi_{\gamma}}$。不同的 Fourier 分量（扇区）在 $Y$ 上沿不同方向以不同的速度移动。因此图景比 Morse 理论的情形更为复杂。

可以证明，生成元 $\Delta=\frac{d}{dt}\big|_{t=0}\varphi^t$ 是 $\Hom_{\mathrm{Dolb}}(E_1,E_2)$ 上的二阶微分算子。当 $g_Y$ 是平坦度量且 $f_1=f_2=0$ 时，可以求出 $\Delta$ 的如下显式公式：
\[
\Delta=i\sum_j\frac{\p^2}{\p x_j\,\p y_j}-\frac{1}{\eps}\sum_j\frac{\p^2}{\p x_j^2}.
\]
看来可以合理地期望：半群 $\varphi^t=e^{t\Delta}$，$t\ge 0$，可以从 $\underline{\Hom}^{\mathrm{alg}}_{\A(Y)}(E_1,E_2)$ 延拓到全部态射空间 $\Hom_{\A(Y)}(E_1,E_2)$。注意 $\Delta$ 不是自伴的，其实部也不是椭圆的。尽管如此，我们期待半群算子 $\varphi^t$ 当 $t\to+\infty$ 时收敛到 Morse 理论情形中的「投影算子」。

\begin{yizhu}[万有丛与平凡化]
$P$ 的定义方式是：先在纤维方向的万有覆盖 $\tilde P$ 上平凡化，再按甲板变换（deck transformation）的乘子下降。这里 $U(1)$-局部系的模空间与对偶环面 $(p^{\vee})^{-1}(y)$ 的等同，正是 T 对偶在纤维上的体现。
\end{yizhu}

\begin{yizhu}[$(0,2)$-部分为何为零]
证明中 $\operatorname{curv}(\nabla_E)^{(0,2)}$ 的系数是 Hessian $\p^2f/\p y_j\p y_k$ 与反对称的 $d\bar z_j\,d\bar z_k$ 的缩并：前者关于 $j,k$ 对称，后者反对称，故和为零。原文把这个常数留作 $\text{const}$，其具体值为 $\pm\pi/\eps$ 量级，与结论无关。
\end{yizhu}

\begin{yizhu}[与 \S8 的对应]
本节的「扇区」与非阿基米德情形中的「单项式 $z_1^{i_1}\dots z_n^{i_n}$」一一对应；「路径的同伦类」对应「Fourier 级数的被加项」。把 $\dbar$ 换成 $d$、把 Dolbeault 复形换成 de Rham 复形之后，几何图像于是完全平行。
\end{yizhu}
